\newpage
\pagestyle{fancy}
\fancyhf{}
\fancyhead[R]{\thepage}
\chapter{PENDAHULUAN} \label{Bab I}

\section{Latar Belakang} \label{I.Latar Belakang}
Teman Tuli atau difabel pendengaran menggunakan bahasa isyarat sebagai alat komunikasi.
Dua bahasa tersebut adalah Sistem Bahasa Isyarat Indonesia (SIBI) dan Bahasa Isyarat Indonesia (BISINDO). 
SIBI lebih umum digunakan pada kegiatan formal dan akademis karena dibuat oleh pemerintah dan Teman Dengar dengan mengadaptasi sistem bahasa isyarat Amerika (\textit{American Sign Language} (ASL)). 
Sementara itu, BISINDO lebih umum digunakan untuk percakapan sehari-hari karena merupakan bahasa yang lahir secara alami dari komunitas tuli.
Namun, keterbatasan pemahaman BISINDO masih menjadi masalah bagi masyarakat umum dalam berkomunikasi dengan Teman Tuli, sehingga dibutuhkan pemanfaatan teknologi untuk mendukung proses penerjemahan bahasa isyarat secara efektif.

% Deep learning sudah luas digunakan
Perkembangan teknologi \textit{computer vision} dan kecerdasan buatan telah banyak dieksplorasi oleh penelitian-penelitian sebelumnya untuk mereka cipta sistem penerjemahan bahasa isyarat ke teks atau suara.
Berbagai studi menunjukkan bahwa model \textit{deep learning} mampu mengekstraksi pola kompleks dari data visual maupun data berurutan seperti citra dan video bahasa isyarat, sehingga menjadi pendekatan dominan dalam pengenalan bahasa isyarat modern.
Penelitian yang dilakukan oleh Meisya memanfaatkan 650 gambar untuk klasifikasi 26 abjad bahasa Indonesia menggunakan XCeption  %masalahanya ini deep learning
Salah satu penelitian terbaru yang diterbitkan oleh \textit{International Journal of Computer Applications} melakukan komparasi terhadap dua model deep learning yaitu \textit{Long Short Term Memory} (LSTM) dan 3D CNN \textit{Convolutional Neural Network} dengan mendapati bahwa 3D CNN mencapai akurasi 92.4\% pada dataset video yang diekstrak landmark \cite{anturkar2025lstm3dcnn}.
Penelitian lain 
% TEkenin pemanfaatan landmark bilang kalaupun banyak yang make gambar sebenernya landmark bisa lebih maju karena udah lebih op daripada gambar saja

% Dominasi pendekatan temporal

% Celah penelitian berbasis spasial, bilang bahwa paper di luar itu make actionr recognition tapi gaada yang kasus BISINDO

% Arah kontribusi penelitian
Mayoritas penelitian yang memanfaatkan \textit{deep learning} untuk membuat penerjemah bahasa isyarat masih menggunakan model \textit{deep learning} berbasis temporal. Hal ini wajar dikarenakan bentuk bahasa isyarat yang selain berurutan juga terkemas dengan baik dalam dimensi waktu. Belum banyak yang memanfaatkan model berbasis spasial seperti CNN atau ResNet dalam klasifikasi gerakan. Padahal, penelitian oleh X memberikan deskripsi menarik tentang potensi penggunaan model berbasis spasial untuk klasifikasi dataset temporal.

Hal ini menarik untuk dipelajari terutama karena tidak ada penelitian terbaru yang mencoba menerjemahkan bahasa isyarat seperti BISINDO menggunakan metode spasial. Perkembangan computer vision juga memungkinkan adanya klasifikasi berbasis landmark mediapipe yang terdiri dari titik-titik koordinat yang secara instan merepresentasikan tubuh manusia seperti wajah, tangan, dan badan. Penelitian ini bertujuan untuk menguji model spasial berbasis neural network untuk klasifikasi gerakan bahasa isyarat menggunakan landmark mediapipe. Diharapkan hasil dari penelitian ini dapat memberikan kontribusi baru dalam dunia penerjemahan bahasa isyarat terutama BISINDO dan juga memberikan gagasan tentang pemanfaatan model spasial dalam klasifikasi atau pengenalan dataset berurutan. 

\section{Rumusan Masalah} \label{I.Rumusan Masalah}
Pemanfaatan model \textit{deep learning} dalam penerjemahan bahasa isyarat seperti BISINDO masih didominasi oleh pendekatan berbasis temporal. 
Sementara itu, pendekatan berbasis spasial untuk penerjemahan bahasa isyarat seperti ResNet masih relatif terbatas.
Belum banyak penelitian yang mengkaji sejauh mana model spasial ResNet-2D mampu menangkap pola gerakan bahasa isyarat ketika representasi data disusun dalam bentuk struktur dua dimensi dari landmark tubuh manusia. 
Oleh karena itu, diperlukan penelitian yang mampu menggali potensi penggunaan model ResNet-2D dalam klasifikasi dan penerjemahan bahasa isyarat Indonesia (BISINDO) ke teks, serta menganalisis pengaruh variasi input landmark terhadap performa model.

% Berdasarkan latar belakang yang telah diuraikan di atas, maka permasalahan penelitian dirumuskan sebagai berikut: \par

% \begin{enumerate}[noitemsep]
% 	\item Apakah model berbasis spasial seperti ResNet-2D mampu menangkap pola gerakan BISINDO ketika representasi data disusun dalam bentuk dua dimensi dari landmark tubuh manusia?
% 	\item Apa pengaruh variasi input landmark (pose, hands, pose+hands, dan full) terhadap performa model ResNet-2D?
% \end{enumerate}


\section{Tujuan Penelitian} \label{I.Tujuan}
Berdasarkan rumusan masalah yang telah diuraikan di atas, maka tujuan dari penelitian ini adalah: \par

\begin{enumerate}[noitemsep]
	\item Mengimplementasikan model spasial berbasis Residual Network (ResNet-2D) untuk klasifikasi gerakan BISINDO.
	\item Menganalisis kinerja model ResNet-2D dalam memanfaatkan representasi spasial landmark MediaPipe pada pengenalan gerakan bahasa isyarat.
	\item Mengkaji pengaruh kombinasi landmark MediaPipe (pose, tangan, wajah) terhadap performa model ResNet-2D.
	\item Mengevaluasi performa model spasial yang digunakan menggunakan metrik evaluasi seperti \textit{Accuarcy, Precision, Recall, F1-Score}. dan \textit{Confussion Matrix}.
\end{enumerate}


\section{Batasan Masalah} \label{I.Batasan}
Adapun batasan masalah dari penelitian ini agar sesuai dengan yang diharapkan adalah sebagai berikut: \par

\begin{enumerate}[noitemsep]
    \item Bahasa isyarat yang digunakan adalah Bahasa Isyarat Indonesia (BISINDO) dengan lingkup cakupan wilayah Jawa Barat.
    \item Jumlah kelas kata yang dijadikan label dibatasi pada 10 kata BISINDO.
    \item Data gerakan bahasa isyarat direpresentasikan menggunakan landmark tubuh yang diekstraksi dengan MediaPipe, meliputi landmark wajah, tangan, dan pose tubuh.
    \item Model utama yang digunakan adalah model spasial berbasis \textit{Residual Network} (ResNet-2D)
    \item Model \textit{Convolutional Neural Network} (CNN) digunakan sebagai pendekatan baseline dan bukan fokus utama dalam analisis komparatif.
    \item Penelitian ini tidak melakukan perbandingan dengan model berbasis temporal seperti LSTM, GRU, atau Transformer.
    \item Penelitian ini tidak melakukan eksplorasi menyeluruh terhadap kombinasi hyperparameter. Konfigurasi hyperparameter yang digunakan relatif seragam untuk menjaga keadilan perbandingan antar model.
    \item Evaluasi performa model difokuskan pada metrik yang umum digunakan dalam tugas klasifikasi seperti \textit{Accuarcy, Precision, Recall, F1-Score}. dan \textit{Confussion Matrix}.
\end{enumerate}

\section{Manfaat Penelitian} \label{I.Manfaat}
Adapun manfaat yang diperoleh dari hasil penelitian ini adalah sebagai berikut: \par

\begin{enumerate}[noitemsep]
    \item Memberikan kontribusi pengetahuan terkait pemanfaatan model spasial berbasis \textit{Residual Network} dalam klasifikasi gerakan BISINDO berbasis landmark MediaPipe.
    \item Memberikan gambaran empiris mengenai pengaruh kombinasi landmark tubuh (pose, tangan, dan wajah) terhadap performa model spasial dalam pengenalan bahasa isyarat.
    \item Menjadi referensi bagi penelitian selanjutnya yang ingin mengeksplorasi pendekatan spasial alam pengolahan data gerakan \textit{action recognition} atau data berurutan tanpa ketergantungan penuh terhadap model berbasis temporal.
    \item  Menjadi bahan evaluasi dan pembelajaran dalam penerapan model deep learning untuk pengenalan pola gerakan berbasis titik koordinat.
\end{enumerate}


\section{Sistematika Penulisan} \label{I.Sistematika}
Sistematika penulisan berisi pembahasan apa yang akan ditulis disetiap Bab. Sistematika pada umumnya berupa paragraf yang setiap paragraf mencerminkan bahasan setiap Bab. \par

\noindent\textbf{Bab I}

Bab ini berisikan penjelasan latar belakang dari topik penelitian yang berlangsung, rumusan masalah dari masalah yang dihadapi pada penjelasan di latar belakang, tujuan dari penelitian, batasan dari penelitian, manfaat dari hasil penelitian, dan sistematika penulisan tugas akhir. \par

\noindent\textbf{Bab II}

Bab ini membahas mengenai teori-teori dan penelitian yang berkaitan dengan penelitian ini.

\noindent\textbf{Bab III}

Bab ini berisikan penjelasan alur kerja sistem, alat dan data yang digunakan, metode yang digunakan, dan rancangan pengujian.

\noindent\textbf{Bab IV}

Bab ini membahas hasil implementasi dan pengujian dari penelitian yang dilakukan.

\noindent\textbf{Bab V}

Bab ini membahas kesimpulan dari hasil penelitian dan juga saran untuk penelitian selanjutnya.